\section{A two-row reduction from 243 to 242 working states}

All labels in this section are inherited labels, rather than consecutive
canonical state numbers.  Let $M$ be the exact combined243 machine at commit
\texttt{e58338db7545909907e07a7603e8eb88625c527a}.  Delete state $15$ and make
the following two changes to retained rows:
\[
\begin{array}{c|c|c|c}
\text{state}&\text{read}&\text{old action}&\text{new action}\\\hline
11&0&(0,R,15)&(0,R,27)\\
27&1&(1,R,39)&(1,R,22).
\end{array}
\]
The two deleted rows are $15/0=(0,R,21)$ and $15/1=(1,R,22)$.
Every other row is unchanged.  Thus the new machine $N$ has exactly $242$
working states, $484$ total binary rows, initial state $95$, and the same
separate halt target $254$.  No row targets deleted state $15$.

The saving reuses the unused one branch of state $27$.  Its zero branch
returns to an earlier tape bit and then resumes forward decoding.  Every
entry that formerly used state $15$ supplies the required earlier zero.
The result is a table transformation, without changing the proof-search
source, register permutation, tape representation, or bootstrap program.

\subsection{The invariant}

For an old-machine configuration $(q,h,T)$ define $I$ by
\[
q\in\{15,18,47,65\}\ \Longrightarrow\ T(h-1)=0,
\qquad
q=27\ \Longrightarrow\ T(h)=0.
\]
This holds at state $95$ for every head position and every tape.  The complete
incoming-row inventory for the constrained states is
\[
\begin{array}{c|l}
\text{target}&\text{incoming rows}\\\hline
15&11/0=(0,R,15)\\
18&12/0=(0,R,18),\quad148/0=(0,R,18)\\
47&33/0=(0,R,47),\quad144/0=(0,R,47),\quad167/0=(0,R,47)\\
65&51/0=(0,R,65),\quad81/0=(0,R,65)\\
27&18/0=(0,L,27),\quad47/0=(0,L,27),\quad65/0=(0,L,27).
\end{array}
\]
An entry into any of $15,18,47,65$ writes zero and moves right, which proves
the required left-zero assertion directly.  An entry into $27$ moves left
from one of $18,47,65$, whose cell immediately left is zero by $I$.  The
write occurs at the departed cell and cannot change that zero.  Every other
target imposes no condition.  Consequently $I$ is preserved by every old
transition.  This is a finite syntactic induction, valid for arbitrary
initial tapes at state $95$; it uses no source-program reachability claim.

\subsection{The two exceptional blocks}

Use all old states except $15$ as comparison checkpoints, with complete
equality of head position and the entire tape.  An old checkpoint at state
$27$ can read only zero by $I$, and the zero transition of $27$ is unchanged.
The only admissible retained row needing a longer comparison is $11/0$.

Suppose this row is entered at head position $h$, so $T(h)=0$.
If $T(h+1)=1$, both machines take two transitions and finish at state $22$
with head $h+2$, preserving the entire tape:
\[
M:\quad11\longrightarrow15\longrightarrow22,
\qquad
N:\quad11\longrightarrow27\longrightarrow22.
\]
If $T(h+1)=0$, the old machine takes two transitions
\[
M:\quad(11,h)\longrightarrow(15,h+1)\longrightarrow(21,h+2),
\]
and the new machine takes four:
\[
N:\quad(11,h)\longrightarrow(27,h+1)\longrightarrow(32,h)
\longrightarrow(45,h+1)\longrightarrow(21,h+2).
\]
The rows used by the last three new transitions are
$27/0=(0,L,32)$, $32/0=(0,R,45)$, and $45/0=(0,R,21)$.
Only cells $h,h+1$ are read or written.  Both are zero in this case and
remain zero, so the whole tape and the endpoint configuration agree for
every assignment to every other cell.  Both clocks are strictly positive.

Every other admissible row out of a checkpoint has exactly the same action
in the two tables and cannot enter deleted state $15$.  The changed row
$27/1$ is inadmissible on the old side by $I$.  The only old row targeting
$15$ is precisely the exceptional row $11/0$.  These assertions are checked
against every row, not inferred from a sample execution.

\subsection{Global theorem}

\paragraph{Theorem.}
For every binary tape $T:\mathbb Z\to\{0,1\}$ and every head $h$, $M$
halts from $(95,h,T)$ if and only if $N$ halts from $(95,h,T)$.

\paragraph{Proof.}
Start with the equal checkpoint configurations in state $95$.  They satisfy
$I$ on the old side.  At an ordinary admissible row use its equal one-step
transitions.  At $11/0$ use the complete two-sided block selected by
$T(h+1)$.  Each block finishes at an equal checkpoint, and the old invariant
holds there because it holds after each old literal transition.  Induction
therefore constructs matching checkpoint computations of both machines.
Every block is finite; the old clock is either one or two, and the new
clock is one, two, or four.  Hence an infinite checkpoint computation
accounts for infinitely many steps on both sides.  A finite halting
computation can only finish at the common halt checkpoint, because no
exceptional block enters the halt state in its interior.  Thus one machine
halts exactly when the other does.\hfill$\square$

The argument is elementary and can be formalized using finite sequences of
configurations and induction on natural numbers.  In particular its usual
arithmetization is provable in PA.  This report supplies an ordinary proof
and a Python checker of its finite premises, not an already kernel-checked
Lean or PA proof artifact.  Composing the theorem with the inherited
ordinary combined243 calibration gives
\[
\mathrm{ZFC}\vdash
\bigl(\neg\operatorname{Halt}(N)\leftrightarrow
       \operatorname{Con}(\mathrm{ZFC})\bigr).
\]
No consistency assumption is needed for the equivalence theorem.  Actual
nonhalting is conditional on consistency, and no absolute consistency or
two-sided independence claim is asserted here.

\subsection{Exact time and space relation}

For any matching finite computation ending at a checkpoint, let $n$ be the
old number of transitions and let $b$ count its visits to $11/0$ for which
the following tape cell is zero.  The new number is exactly $n+2b$.
Every other comparison block has equal old and new lengths.  In particular,
if either machine halts, their halting times $n,n'$ satisfy
\[
n'=n+2b,\qquad n\leq n'\leq2n.
\]
The longer block visits no new position: both machines visit precisely
$\{h,h+1,h+2\}$ in that block.  Consequently the sets of head positions
visited up to matching checkpoints are equal, as are their finite minimum
and maximum positions when these exist.

For the prescribed blank bootstrap, the old computation has exactly six
longer-block entries.  Independent execution gives $407\,967$ old and
$407\,979$ new transitions; both reach the same full root frame, with
$351$ one-cells, head $-5$, and physical registers
$(0,0,0,3,0,0,3,0,0)$.  Both traces visit head positions from $-15$ to $701$.
The numerical relation agrees with $407\,979=407\,967+2\cdot6$.

\subsection{Reproducible evidence}

The accompanying \texttt{verify\_return242.py} checks exact table hashes,
totality, the full incoming inventory proving $I$, both universal finite
block identities, and every ordinary-row condition.  In its symbolic
calculation, an unconstrained tape cell may never be read; unmentioned cells
are preserved by construction.  Thus those local checks quantify over all
completions of the partial tape, rather than testing a bounded collection
of complete input tapes.  A separate sparse blank replay verifies that the
old and new bootstraps reach precisely equal prescribed root-frame configurations.
The blank replay is supplementary finite evidence and is not a replacement
for the global induction proof.
